% GENERATED FILE. Edit canonical lessons, then run npm run generate:books.
% canonical-source-hash: fnv1a64:ee05cc9a04645095
% canonical-lessons: RU-C105-flour, RU-C105-chocolate, RU-C105-biscuits, RU-C105-sandwich, RU-C105-sauce

\chapter{Flour, Chocolate, Biscuits, Sandwich, Sauce}
\label{ch:ru-flour-chocolate-biscuits-sandwich-sauce}
\hlchaptermodality{105}

\begin{chapteropening}
\textbf{By the end of this chapter:} \emph{I can say flour, chocolate, biscuits, a sandwich and a sauce.}

Complete the last lesson of chapter 105: I can say flour, chocolate, biscuits, a sandwich and a sauce.
\end{chapteropening}

\section[\texorpdfstring{\ru{мука} (muká)}{мука (muká)}]{\ru{мука} (muká) — flour}

\label{lesson:RU-C105-flour}



\begin{quote}
Before the new one: say the Russian for a melon, then the Russian for garlic.
\end{quote}

\subsection*{You'll want to know: \ru{мука}}
\textbf{\ru{мука}} (\emph{muká}) — \textquotedblleft{}flour\textquotedblright{}.

\begin{cousinweb}
\textbf{\ru{мука}}, \textquotedblleft{}flour\textquotedblright{}, stressed on the last syllable $\leftarrow$ Proto-Slavic \emph{*mǫka} — inherited, from the root of \emph{\ru{мягкий}}, \textquotedblleft{}soft\textquotedblright{}. Vasmer relates it to Lithuanian \emph{minkyti}, \textquotedblleft{}to knead\textquotedblright{}, and Greek \emph{mássō}, \textquotedblleft{}I knead\textquotedblright{}. \emph{\ru{мука}}, stressed on the first syllable, \textquotedblleft{}torment\textquotedblright{}, is a related word.
\end{cousinweb}

\begin{grammarlens}[title={What you've built}]
One more word for this chapter.
\end{grammarlens}

\subsection*{Your turn}
\begin{itemize}
\raggedright
  \item \emph{Say it:} \emph{muká}
  \item \emph{Say it:} \emph{muká}, once more
  \item \emph{Say it:} say \emph{muká}
  \item \emph{Recall:} say the Russian for a melon, then the Russian for garlic, then say \emph{muká} again
\end{itemize}

\begin{tcolorbox}[breakable,colback=teal!4,colframe=teal!35!black,title={Before you move on}]
What does \ru{мука} mean? (\textquotedblleft{}Flour\textquotedblright{}.) Say it once more.
\end{tcolorbox}
\section[\texorpdfstring{\ru{шоколад} (shokolád)}{шоколад (shokolád)}]{\ru{шоколад} (shokolád) — chocolate}

\label{lesson:RU-C105-chocolate}



\begin{quote}
Before the new one: say the Russian for garlic, then the Russian for flour.
\end{quote}

\subsection*{You'll want to know: \ru{шоколад}}
\textbf{\ru{шоколад}} (\emph{shokolád}) — \textquotedblleft{}chocolate\textquotedblright{}.

\begin{cousinweb}
\textbf{\ru{шоколад}} — from Peter the Great's time. The \emph{\ru{ш}-} form came through German \emph{Schokolade}, the old \emph{\ru{ч}-} form from Italian \emph{cioccolata}; both go back to Spanish \emph{chocolate}, from Nahuatl, the language of the Aztecs (Vasmer).
\end{cousinweb}

\begin{grammarlens}[title={What you've built}]
One more word for this chapter.
\end{grammarlens}

\subsection*{Your turn}
\begin{itemize}
\raggedright
  \item \emph{Say it:} \emph{shokolád}
  \item \emph{Say it:} \emph{shokolád}, once more
  \item \emph{Say it:} say \emph{shokolád}
  \item \emph{Recall:} say the Russian for garlic, then the Russian for flour, then say \emph{shokolád} again
\end{itemize}

\begin{tcolorbox}[breakable,colback=teal!4,colframe=teal!35!black,title={Before you move on}]
What does \ru{шоколад} mean? (\textquotedblleft{}Chocolate\textquotedblright{}.) Say it once more.
\end{tcolorbox}
\section[\texorpdfstring{\ru{печенье} (pechénye)}{печенье (pechénye)}]{\ru{печенье} (pechénye) — biscuits}

\label{lesson:RU-C105-biscuits}



\begin{quote}
Before the new one: say the Russian for flour, then the Russian for chocolate.
\end{quote}

\subsection*{You'll want to know: \ru{печенье}}
\textbf{\ru{печенье}} (\emph{pechénye}) — \textquotedblleft{}biscuits\textquotedblright{}.

\begin{cousinweb}
\textbf{\ru{печенье}} is a noun from \emph{\ru{печь}}, \textquotedblleft{}to bake\textquotedblright{}: \textquotedblleft{}something baked\textquotedblright{}. \emph{\ru{печь}} $\leftarrow$ Proto-Slavic \emph{*pekti}, which Vasmer relates to Sanskrit \emph{paktí-}, \textquotedblleft{}cooking\textquotedblright{}, and Greek \emph{pépsis}, \textquotedblleft{}cooking\textquotedblright{}.
\end{cousinweb}

\begin{grammarlens}[title={What you've built}]
One more word for this chapter.
\end{grammarlens}

\subsection*{Your turn}
\begin{itemize}
\raggedright
  \item \emph{Say it:} \emph{pechénye}
  \item \emph{Say it:} \emph{pechénye}, once more
  \item \emph{Say it:} say \emph{pechénye}
  \item \emph{Recall:} say the Russian for flour, then the Russian for chocolate, then say \emph{pechénye} again
\end{itemize}

\begin{tcolorbox}[breakable,colback=teal!4,colframe=teal!35!black,title={Before you move on}]
What does \ru{печенье} mean? (\textquotedblleft{}Biscuits\textquotedblright{}.) Say it once more.
\end{tcolorbox}
\section[\texorpdfstring{\ru{бутерброд} (buterbród)}{бутерброд (buterbród)}]{\ru{бутерброд} (buterbród) — a sandwich}

\label{lesson:RU-C105-sandwich}



\begin{quote}
Before the new one: say the Russian for chocolate, then the Russian for biscuits.
\end{quote}

\subsection*{You'll want to know: \ru{бутерброд}}
\textbf{\ru{бутерброд}} (\emph{buterbród}) — \textquotedblleft{}a sandwich\textquotedblright{}.

\begin{cousinweb}
\textbf{\ru{бутерброд}} comes from German \emph{Butterbrot}, \textquotedblleft{}buttered bread\textquotedblright{} (Vasmer): \emph{Butter}, \textquotedblleft{}butter\textquotedblright{}, with \emph{Brot}, \textquotedblleft{}bread\textquotedblright{}.
\end{cousinweb}

\begin{grammarlens}[title={What you've built}]
One more word for this chapter.
\end{grammarlens}

\subsection*{Your turn}
\begin{itemize}
\raggedright
  \item \emph{Say it:} \emph{buterbród}
  \item \emph{Say it:} \emph{buterbród}, once more
  \item \emph{Say it:} say \emph{buterbród}
  \item \emph{Recall:} say the Russian for chocolate, then the Russian for biscuits, then say \emph{buterbród} again
\end{itemize}

\begin{tcolorbox}[breakable,colback=teal!4,colframe=teal!35!black,title={Before you move on}]
What does \ru{бутерброд} mean? (\textquotedblleft{}A sandwich\textquotedblright{}.) Say it once more.
\end{tcolorbox}
\section[\texorpdfstring{\ru{соус} (sóus)}{соус (sóus)}]{\ru{соус} (sóus) — a sauce}

\label{lesson:RU-C105-sauce}



\begin{quote}
Before the new one: say the Russian for biscuits, then the Russian for a sandwich.
\end{quote}

\subsection*{You'll want to know: \ru{соус}}
\textbf{\ru{соус}} (\emph{sóus}) — \textquotedblleft{}a sauce\textquotedblright{}.

\begin{cousinweb}
\textbf{\ru{соус}} comes from French \emph{sauce}, from Latin \emph{salsa}, \textquotedblleft{}salted (broth)\textquotedblright{}, perhaps through English \emph{sauce} (Vasmer).
\end{cousinweb}

\begin{grammarlens}[title={What you've built}]
One more word for this chapter.
\end{grammarlens}

\subsection*{Your turn}
\begin{itemize}
\raggedright
  \item \emph{Say it:} \emph{sóus}
  \item \emph{Say it:} \emph{sóus}, once more
  \item \emph{Say it:} say \emph{sóus}
  \item \emph{Recall:} say the Russian for biscuits, then the Russian for a sandwich, then say \emph{sóus} again
\end{itemize}

\begin{tcolorbox}[breakable,colback=teal!4,colframe=teal!35!black,title={Before you move on}]
What does \ru{соус} mean? (\textquotedblleft{}A sauce\textquotedblright{}.) Say it once more.
\end{tcolorbox}
